\documentclass[10pt]{article}
\providecommand{\FIGDIR}{fig}
\newcommand{\DIGESTDATE}{22 August 2026}
% ===================== PM Research Watch - shared preamble =====================
\usepackage[T1]{fontenc}
\usepackage[utf8]{inputenc}
\usepackage{amsmath,amssymb}
\usepackage{tgpagella}
\usepackage{tgheros}
\renewcommand{\sfdefault}{qhv}
\renewcommand{\ttdefault}{lmtt}
\usepackage{microtype}
\usepackage[english]{babel}

\usepackage[a4paper,top=17mm,bottom=17mm,left=17mm,right=17mm,
            headsep=6mm,footskip=10mm]{geometry}
\usepackage{graphicx}
\usepackage{xcolor}
\usepackage{colortbl}
\usepackage{tikz}
\usetikzlibrary{calc}
\usepackage[most]{tcolorbox}
\usepackage{titlesec}
\usepackage{fancyhdr}
\usepackage{booktabs}
\usepackage{array}
\usepackage{longtable}
\usepackage{enumitem}
\usepackage{multicol}
\usepackage{ragged2e}
\usepackage[hidelinks]{hyperref}

% ------------------------------------------------------------------ palette
\definecolor{Ink}    {HTML}{14262E}
\definecolor{Deep}   {HTML}{0F4C5C}
\definecolor{Teal}   {HTML}{1F7A8C}
\definecolor{Sky}    {HTML}{4FA3C4}
\definecolor{Sage}   {HTML}{6FA287}
\definecolor{Amber}  {HTML}{E3A008}
\definecolor{Clay}   {HTML}{B5643C}
\definecolor{Coral}  {HTML}{D1495B}
\definecolor{Violet} {HTML}{6A5B9E}
\definecolor{Slate}  {HTML}{4A6572}
\definecolor{Paper}  {HTML}{FBF9F5}
\definecolor{Mist}   {HTML}{EEEAE2}
\definecolor{Rule}   {HTML}{D9D3C9}

\pagecolor{Paper}
\color{Ink}

\hypersetup{colorlinks=true, linkcolor=Deep, urlcolor=Teal, citecolor=Teal,
            pdftitle={PM Research Watch}, pdfauthor={Automated daily digest}}

% ------------------------------------------------------------------ headings
\titleformat{\section}
  {\sffamily\bfseries\large\color{Deep}}{}{0pt}
  {}[{\vspace{1.2mm}\color{Amber}\titlerule[1.1pt]}]
\titlespacing*{\section}{0pt}{5.5mm}{2.8mm}

\titleformat{\subsection}
  {\sffamily\bfseries\normalsize\color{Teal}}{}{0pt}{}
\titlespacing*{\subsection}{0pt}{3.6mm}{1.4mm}

\setlist[itemize]{leftmargin=4.2mm, itemsep=1.1mm, topsep=1.3mm, parsep=0pt,
                  label=\textcolor{Teal}{\footnotesize$\blacktriangleright$}}

% ------------------------------------------------------------------ header/footer
\pagestyle{fancy}
\fancyhf{}
\renewcommand{\headrulewidth}{0pt}
\renewcommand{\footrulewidth}{0pt}
\fancyhead[L]{\sffamily\scriptsize\color{Slate}PM RESEARCH WATCH\;\textcolor{Rule}{\textbar}\;\DIGESTDATE}
\fancyhead[R]{\sffamily\scriptsize\color{Slate}Particulate matter monitoring \& health effects}
\fancyfoot[C]{\sffamily\scriptsize\color{Slate}\thepage}

% ------------------------------------------------------------------ components
% metric tile
\newcommand{\metric}[3]{%
  \begin{tcolorbox}[enhanced, colback=white, colframe=Rule, boxrule=0.5pt,
    arc=1.6mm, left=2.2mm, right=2.2mm, top=1.8mm, bottom=1.8mm,
    borderline west={1.7mm}{0pt}{#3}, nobeforeafter, valign=center,
    height=17mm, width=\linewidth]
    \RaggedRight
    {\sffamily\bfseries\fontsize{17}{17}\selectfont\color{#3}#1}\\[0.5mm]
    {\sffamily\fontsize{6.5}{7.4}\selectfont\color{Slate}#2}
  \end{tcolorbox}}

% pull-quote / key insight box
\newtcolorbox{keybox}[1][]{enhanced, colback=Mist, colframe=Mist,
  arc=1.6mm, left=3.2mm, right=3.2mm, top=2.4mm, bottom=2.4mm,
  borderline west={1.7mm}{0pt}{Amber}, #1}

% subtopic band
\newcommand{\band}[2]{%
  \vspace{3.2mm}\noindent
  \begin{tcolorbox}[enhanced, colback=#2, colframe=#2, arc=1.2mm,
    left=2.6mm, right=2.6mm, top=1.1mm, bottom=1.1mm, boxrule=0pt,
    width=\linewidth, nobeforeafter]
    {\sffamily\bfseries\small\color{white}#1}
  \end{tcolorbox}\vspace{1.6mm}}

% one paper entry: {short cite}{title}{journal}{doi}{body}
\newcommand{\paperentry}[5]{%
  \par\vspace{1.9mm}\noindent
  {\sffamily\bfseries\footnotesize\color{Ink}#1}\;%
  {\sffamily\fontsize{7}{8}\selectfont\color{Slate}#3}\\[0.4mm]
  {\footnotesize\itshape\color{Deep}#2}\\[0.6mm]
  {\footnotesize #4}\\[0.3mm]
  {\fontsize{6.6}{7.6}\selectfont\color{Slate}\texttt{doi:}\,\href{https://doi.org/#5}{\texttt{\detokenize{#5}}}}
  \par}

% caption style
\newcommand{\figcap}[1]{\par\vspace{1.0mm}
  {\fontsize{7.2}{8.6}\selectfont\color{Slate}\textsf{\textbf{Fig.}}\; #1\par}\vspace{1.5mm}}

\providecommand{\permil}{\textperthousand}

\setlength{\parindent}{0pt}
\setlength{\parskip}{0pt}
\renewcommand{\arraystretch}{1.18}


\begin{document}

% ==================================================================== MASTHEAD
\thispagestyle{empty}
\begin{tikzpicture}[remember picture, overlay]
  \fill[Deep] (current page.north west) rectangle
      ($(current page.north east)+(0,-67mm)$);
  \fill[Amber] ($(current page.north west)+(0,-67mm)$) rectangle
      ($(current page.north east)+(0,-68.4mm)$);
  \foreach \x/\y/\r in {22/12/0.9, 41/26/1.5, 63/9/0.7, 88/31/1.1, 112/17/1.9,
                        138/29/0.8, 159/13/1.3, 178/54/1.0, 196/21/1.6,
                        30/44/0.6, 74/57/1.2, 125/48/0.9, 168/6/0.8, 191/59/0.7,
                        52/18/1.0, 100/60/0.6, 148/20/1.1, 66/49/1.4, 205/45/0.8}
    \fill[Sky, opacity=0.30] ($(current page.north west)+(\x mm,-\y mm)$) circle (\r mm);
\end{tikzpicture}

\vspace*{4mm}
{\sffamily\fontsize{7.6}{9}\selectfont\color{Sky}%
 AUTOMATED DAILY LITERATURE DIGEST\;\textbullet\;PhD RESEARCH WATCH\par}
\vspace{2.2mm}
{\sffamily\bfseries\fontsize{27}{29}\selectfont\color{white}%
 Particulate Matter\par}
{\sffamily\fontsize{27}{29}\selectfont\color{Sky}%
 Monitoring \& Health Effects\par}
\vspace{5.0mm}
{\sffamily\fontsize{8.4}{10}\selectfont\color{white}%
 \mbox{Issue 2026-08-22}\;\;\textcolor{Amber}{\textbar}\;\;\mbox{Entry window 22 Aug 2026}\;\;
 \textcolor{Amber}{\textbar}\;\;\mbox{13 in scope, 44 screened out}\par}
\vspace{18mm}

% -------------------------------------------------------------- metric strip
\noindent
\begin{minipage}[t]{0.190\linewidth}\metric{13}{RECORDS IN SCOPE\\(44 SCREENED OUT)}{Deep}\end{minipage}\hfill
\begin{minipage}[t]{0.190\linewidth}\metric{7}{OF 10 CLUSTERS\\OCCUPIED}{Teal}\end{minipage}\hfill
\begin{minipage}[t]{0.190\linewidth}\metric{6}{ABSTRACT-FREE\\DEPOSITS (46\%)}{Sky}\end{minipage}\hfill
\begin{minipage}[t]{0.190\linewidth}\metric{5}{RISK ESTIMATES\\ON THE PLOT}{Amber}\end{minipage}\hfill
\begin{minipage}[t]{0.190\linewidth}\metric{3}{INFECTION-RELATED\\RECORDS}{Clay}\end{minipage}

\vspace{6mm}

% -------------------------------------------------------------- signal of day
\section*{\sffamily\bfseries\color{Deep}Signal of the day}
\vspace{-1.5mm}{\color{Amber}\rule{\linewidth}{1.1pt}}\vspace{2.5mm}

\begin{keybox}
{\sffamily\bfseries\small\color{Deep}Three papers, three scales, one claim: particulate
exposure degrades the innate pulmonary defence against infection --- and today one of them
carries a rescue.}\\[1.4mm]
\footnotesize
Zhang et al.\ (\textit{Environ Sci Technol}) followed 3{,}583 hospitalised patients with
severe fever with thrombocytopenia syndrome and found long-term pre-admission PM$_{2.5}$
associated with \textbf{HR 1.16 (95\% CI 1.04--1.30) per 10\,$\mu$g/m$^3$} for
SFTS-attributable death, with viral load partially mediating; the paired mouse and cell
work shows PM$_{2.5}$ pre-exposure raising SFTSV replication and tissue injury, implicates
mitochondrial redox and ROS pathways, and \emph{reverses the effect with
N-acetylcysteine}. Masci et al.\ (\textit{Environ Toxicol Pharmacol}) reach the same
place from nutrition: in rats, chronic 20\% caloric restriction alone reproduced the
alveolar-macrophage phenotype that residual oil fly ash produced in well-nourished
animals --- reduced viability, lower Nrf2, higher TNF$\alpha$ --- and undernourished
macrophages then \textbf{failed to mount a response to a Poly(I:C) viral challenge
regardless of PM exposure}, i.e.\ the two insults are not additive because the first one
has already saturated the damage. Introna et al.\ (\textit{ACS EST Air}) supply the
dosimetric caveat that both need: exposing lung cells to subway-platform PM$_{2.5}$ at a
realistic \emph{in situ} air-liquid-interface dose of 0.12\,$\pm$\,0.07\,$\mu$g/cm$^2$
produced \textbf{no measurable change in viability or inflammatory markers at 24\,h},
while the same particles submerged at laboratory doses moved four cytokines. The
uncomfortable synthesis is that the epidemiology is real, the mechanism is plausible, and
the standard in-vitro dose is nowhere near the ambient one.
\end{keybox}

\vspace{3mm}

% -------------------------------------------------------------- exec brief
\section*{\sffamily\bfseries\color{Deep}What changed today}
\vspace{-1.5mm}{\color{Amber}\rule{\linewidth}{1.1pt}}\vspace{2.5mm}

\footnotesize
\begin{itemize}
\item \textbf{Six of thirteen records arrived without an abstract --- the highest
abstract-free share on record.} All six came through the Crossref by-ISSN sweep
(\textit{Atmos Environ}, \textit{Atmos Res}, \textit{Atmos Pollut Res},
\textit{Environ Pollut}), and Saturday's thin PubMed leg (2 records) meant the sweep
carried 8 of 13. They are shipped as \textit{Metadata only} with design, sample and
endpoint left unextracted rather than inferred. Two of them --- aerosol pH and preterm
birth, and ammonia-driven ROS generation in secondary organic aerosol --- are the day's
most interesting titles and the least characterisable records in it.

\item \textbf{A household-fuel transition with a measurable physiological dividend.}
In CHARLS, moving from solid to clean cooking fuel attenuated peak-expiratory-flow decline
by \textbf{11.0\,L/min (95\% CI 4.7--17.3)} over two years and was associated with
\textbf{OR 0.80 for chronic lung disease and 0.80 for heart disease} over the long term.
This is one of the few PM-relevant exposures where the intervention, the physiological
intermediate and the disease endpoint are all measured in the same cohort.

\item \textbf{An outdoor-activity framing of wildfire smoke, at a two-orders-of-magnitude
ratio.} A preprint on the July 2026 Gironde--Landes fires reports the footprint over the
EAQI ``Poor'' threshold (hourly PM$_{2.5}$ $>$ 50\,$\mu$g/m$^3$) at \textbf{149 times the
burned area}, extending roughly 370\,km downwind. The methodological point is the
correction the author makes against himself: the 10\,km CAMS analysis reproduces the
episode (daily $r = 0.83$) but underestimates peaks by a median factor of $\sim$2.2, so
the mapped heat--smoke relay is a conservative floor, not an estimate.

\item \textbf{Childhood eczema, with traffic proximity separated from the household
cluster.} A 1{,}227-child Porto survey isolates \textbf{aOR 1.32 (1.11--1.48) for
residence within 300\,m of a traffic zone} alongside \textbf{aOR 1.48 (1.26--1.73) for
postnatal secondhand smoke}, after adjustment that also absorbs family history
(aOR 3.53), income and maternal education. Cross-sectional and questionnaire-based ---
but the traffic term surviving that adjustment set is the part worth citing.

\item \textbf{Front-matter contamination in the newly promoted ES\&T feed, caught and
closed.} The first run with the widened ISSN list pulled \textbf{154 works from
\textit{Environ Sci Technol}, of which 142 were mastheads and issue-information stubs}
from volumes 42--52 being re-deposited in bulk. They carry DOIs and titles, so nothing
downstream would have rejected them. A title-pattern filter now drops them at the source.
\end{itemize}

\clearpage

% ==================================================================== ANALYTICS
\section*{\sffamily\bfseries\color{Deep}Corpus analytics}
\vspace{-1.5mm}{\color{Amber}\rule{\linewidth}{1.1pt}}\vspace{3mm}

\noindent\includegraphics[width=\linewidth]{\FIGDIR/f1_subtopics.png}
\figcap{Subtopic assignment of the 13 in-scope records; single-label by dominant
contribution. \textit{Exposure assessment \& modelling} takes 5, and the health side is
unusually evenly spread --- mechanistic toxicology 2, and one record each in
cardiovascular/metabolic, respiratory/allergic, reproductive/developmental and other
clinical endpoints. Seven of ten canonical clusters are occupied on a 13-record day,
which is a wider spread than the 19-record issue that preceded it.}

\vspace{2mm}
\noindent\includegraphics[width=\linewidth]{\FIGDIR/f2_design_endpoint.png}
\figcap{Left --- study architecture. \textit{Metadata only} is the largest single slice at
6 of 13 (46\%), the highest share the watch has recorded. That slice is a statement about
what publishers deposited, not about what was published: every one of the six is a
Crossref journal-sweep record whose publisher supplied a title and DOI but no abstract.
Right --- health endpoint by organ system. Four records carry a respiratory endpoint, and
the \textit{Not extracted} bar is exactly the six abstract-free records.}

\vspace{2mm}
\noindent\includegraphics[width=\linewidth]{\FIGDIR/f3_metric_geography.png}
\figcap{Left --- particle metric. Six records are PM$_{2.5}$-specific and seven are
unspeciated PM or a PM-proxy exposure (household solid fuel, traffic proximity, ROFA,
mineral dust, aerosol acidity). Right --- geographic provenance. Europe leads with 4,
China with 3, and the Middle East / North Africa contributes 2 --- an Iranian
winter-PM$_{2.5}$ synoptic study and a Mesopotamian dust projection, both
metadata-only. The three ``global / multi-region'' records include two whose study
location is genuinely not stated in the deposited metadata.}

\vspace{2mm}
\noindent\includegraphics[width=\linewidth]{\FIGDIR/f5_forest.png}
\figcap{All five forest-plottable estimates, on a common logarithmic axis.
\textbf{The exposure contrasts are not mutually comparable} --- a per-10\,$\mu$g/m$^3$
ambient PM$_{2.5}$ contrast, two binary household/residential exposures, and a
fuel-transition contrast --- so read direction and precision, not relative magnitude.
The two Liu et al.\ odds ratios are protective and their published intervals are given to
one decimal (0.7--1.0), so the upper bound touches unity even though the paper reports
$p = 0.044$ and $p = 0.028$; they are transcribed as published rather than back-computed
from the $p$-values.}

\vspace{2mm}
\noindent\includegraphics[width=\linewidth]{\FIGDIR/f6_lifecourse.png}
\figcap{Life-course windows addressed, counted only where a record's study population is
explicitly anchored at that stage. Four of five windows are occupied --- the widest
life-course coverage since 14 August --- and the day spans the whole arc from an
aerosol-acidity/preterm-birth record through a paediatric eczema survey to a
middle-aged-and-older cohort. Only adolescence is empty.}

\vspace{2mm}
\noindent\includegraphics[width=\linewidth]{\FIGDIR/f4_heatmap.png}
\figcap{Subtopic $\times$ study-architecture cross-tabulation. The concentration of the
\textit{Metadata only} column across four different subtopic rows is the structural
feature of the day: the abstract-free deposits are not clustered in one corner of the
corpus, so the gap they leave is spread across the whole issue rather than confined to
one topic.}

\vspace{3mm}
% ==================================================================== DIGEST
\section*{\sffamily\bfseries\color{Deep}Paper-level digest}
\vspace{-1.5mm}{\color{Amber}\rule{\linewidth}{1.1pt}}\vspace{1mm}

{\fontsize{7.4}{9}\selectfont\color{Slate}Ordered by cluster. Each entry gives the
methodological core and the finding that matters, not an abstract paraphrase.\par}

\band{Mechanistic toxicology\hfill 2 records}{Violet}

\paperentry{Introna et al.\ 2026}
{Novel approach to assess the toxicity of PM$_{2.5}$ on a subway platform using a mobile air-liquid interface system and submerged exposure}
{ACS EST Air}
{A mobile ALI exposure system operated \emph{on the platform itself}, with A549
monoculture and A549 + dTHP-1 coculture, benchmarked against conventional submerged
laboratory exposure to the same metal-rich mechanical-wear particles. At the realistic
in-situ ALI dose of \textbf{0.12 $\pm$ 0.07\,$\mu$g/cm$^2$ there was no change in
viability or inflammatory markers at 24\,h}. Submerged at higher dose, viability fell in
both cell types; A549 released IL-6 and TNF$\alpha$ at 100\,$\mu$g/mL, and dTHP-1
released IL-8 at the lowest dose tested (10\,$\mu$g/mL) with all four cytokines elevated
at 100. The authors' own reading is the right one --- longer in-situ exposures are needed
to reach comparable dose --- and the wider implication is that the submerged literature's
dose range is not the ambient one.}
{10.1021/acsestair.6c00018}

\paperentry{Masci et al.\ 2026}
{Air pollution and undernutrition induce rat alveolar macrophage functional dysregulation and compromise their response to a viral challenge}
{Environ Toxicol Pharmacol}
{Male Wistar rats, control vs.\ nutritional growth retardation (20\% caloric restriction,
4 weeks), each with intranasal residual oil fly ash or vehicle; isolated alveolar
macrophages assayed for viability, TNF$\alpha$, oxidative stress, Nrf2 and LC3, then
challenged with Poly(I:C). Undernutrition alone reproduced what ROFA did to well-nourished
animals. \textbf{ROFA added little or nothing on top of undernutrition} --- an attenuated
rather than additive response --- and after Poly(I:C), control macrophages mounted an
inflammatory response that prior ROFA \emph{amplified}, while undernourished macrophages
failed to respond at all regardless of exposure. The interaction structure, not the main
effect, is the finding: in a population that is already nutritionally compromised, the
marginal immunological cost of PM may be small because the defence is already gone.}
{10.1016/j.etap.2026.105138}

\band{Other clinical endpoints\hfill 1 record}{Slate}

\paperentry{Zhang et al.\ 2026}
{PM$_{2.5}$ exposure and the prognosis of severe fever with thrombocytopenia syndrome}
{Environ Sci Technol}
{A cohort of \textbf{3{,}583 hospitalised SFTS patients} with long-term pre-admission
ambient PM$_{2.5}$: \textbf{HR 1.16 (95\% CI 1.04--1.30) per 10\,$\mu$g/m$^3$} for
SFTS-attributable death, with exploratory mediation implicating viral load. A
counterfactual burden assessment estimates \textbf{137 avoidable SFTS deaths in Xinyang,
2010--2022} had annual PM$_{2.5}$ met the Chinese national standard. The experimental arm
shows PM$_{2.5}$ pre-exposure increasing SFTSV replication and tissue injury;
transcriptomics implicate mitochondrial redox and ROS pathways, and ROS/MDA measurement
plus N-acetylcysteine intervention support oxidative priming as the upstream step.
Emerging-infection prognosis is a genuinely new endpoint class for this watch.}
{10.1021/acs.est.6c03214}

\band{Cardiovascular \& metabolic\hfill 1 record}{Teal}

\paperentry{Liu et al.\ 2026b}
{Household clean cooking fuel transition and attenuated cardiopulmonary function decline in middle-aged and older adults}
{Indoor Air}
{CHARLS 2011--2020, with two independent subcohorts built from baseline solid-fuel users:
a short-term physiological cohort (n = 6{,}534; 27.4\% transitioned by 2013) and a
long-term disease cohort (n = 6{,}362; 33.5\% transitioned by 2015). Transition attenuated
peak-expiratory-flow decline by \textbf{11.0\,L/min (95\% CI 4.7--17.3)} and lowered
systolic blood pressure modestly ($\beta = -1.7$, $p = 0.017$). Long term:
\textbf{OR 0.80 (0.7--1.0, $p = 0.044$) for chronic lung disease} and \textbf{0.80
(0.7--1.0, $p = 0.028$) for heart disease}, with no association for hypertension;
protective effects were larger in non-smokers and non-drinkers. The two-subcohort design
is what lets the physiological intermediate and the disease endpoint be read together.}
{10.1155/ina/8898458}

\band{Respiratory \& allergic\hfill 1 record}{Sage}

\paperentry{Hatem et al.\ 2026}
{Child characteristics, household and home environmental factors associated with eczema}
{Eur J Pediatr}
{1{,}227 children aged 5--11 from 20 primary schools in Porto (46.5\% response), analysed
with mixed-effects logistic regression with a school random intercept and multiple
imputation. Independently associated with higher odds: \textbf{residence within 300\,m of
a traffic zone, aOR 1.32 (1.11--1.48)}; \textbf{postnatal secondhand smoke, aOR 1.48
(1.26--1.73)}; family history 3.53; higher family income 1.40; all-purpose cleaner use
1.16. Lower odds: maternal education 0.68, detached housing 0.67, fur blanket use 0.85.
Household dampness interacted with residence type and winter ventilation frequency.
Cross-sectional and self-reported, so temporality is unavailable --- but the traffic term
surviving a socioeconomic adjustment set this dense is the citable part.}
{10.1007/s00431-026-07327-1}

\band{Reproductive \& developmental\hfill 1 record}{Sky}

\paperentry{Han et al.\ 2026}
{Aerosol pH associated with preterm birth risk: unmasking the hidden role of atmospheric acidity in maternal-fetal health}
{Environ Pollut}
{\textbf{Metadata only.} Surfaced through the Crossref by-ISSN sweep with no abstract
deposited; title, journal, authorship and DOI verified, nothing else extracted. Worth
flagging by title alone: aerosol acidity as an exposure metric in a reproductive endpoint
is not something this watch has seen before, and it pairs directly with Chai et al.'s
pH-governed partitioning result in yesterday's issue. Queued for abstract retry.}
{10.1016/j.envpol.2026.129011}

\clearpage

\band{Sensing, forecasting \& instrumentation\hfill 2 records}{Amber}

\paperentry{Fei et al.\ 2026}
{SHAP-interpretable machine learning integrating exposures and multi-omics in COPD-lung cancer comorbidity}
{Front Immunol}
{A sequential cross-scale framework: GBD epidemiology 1990--2021 across 204 countries,
exposome-wide random-forest classification with SHAP interpretation, immune
characterisation via CIBERSORT and single-cell RNA sequencing, candidate-gene
prioritisation by SMR and LASSO, then wet-lab validation including a macrophage-epithelial
Transwell co-culture. \textbf{SHAP identified smoking, particulate matter pollution and
residential radon as the major shared risk factors}; TREM1 and ODF2L emerged as shared hub
genes, downregulated in both diseases, and macrophage-specific silencing of either
attenuated macrophage-driven A549 migration and proliferation. The exposure side is
GBD-level and coarse --- this is not an exposure-assessment paper --- but it is a clean
worked example of putting an exposome variable through an interpretable-ML pipeline and
ending at a testable target.}
{10.3389/fimmu.2026.1875146}

\paperentry{Zeng et al.\ 2026}
{Aerosol vertical distributions determine the seasonal relationships between column AOD and surface PM$_{2.5}$ over Europe}
{Atmos Environ}
{\textbf{Metadata only.} No abstract deposited. The title states a result that matters
directly to satellite-based PM$_{2.5}$ estimation --- the AOD-to-surface conversion being
governed by vertical structure and therefore seasonally non-stationary --- so this is the
highest-priority record of the six for abstract retry or full-text retrieval on the next
run.}
{10.1016/j.atmosenv.2026.122305}

\band{Exposure assessment \& modelling\hfill 5 records}{Deep}

\paperentry{Defrance 2026 (preprint)}
{Beyond the fire perimeter: the smoke footprint that constrains outdoor activity dwarfs the burned area of the July 2026 Gironde--Landes wildfire episode}
{Research Square}
{ERA5-Land wet-bulb globe temperature combined with CAMS analysis PM$_{2.5}$ at hourly
resolution over 25--31 July 2026, against 40 background monitoring stations. The two
fires burned 340\,km$^2$ and forced $\sim$220{,}000 evacuations; the footprint above the
European AQI ``Poor'' threshold (hourly PM$_{2.5}$ $>$ 50\,$\mu$g/m$^3$, where reduced
outdoor activity is advised) reached \textbf{149 times the burned area} and about
\textbf{370\,km downwind}, with 139$\times$ on the EU 2030 daily limit of 25\,$\mu$g/m$^3$
and at least 121$\times$ against the larger operational fire perimeter. Plume co-location
rose from 73\% hourly to 94\% daily. At the stations, a same-day heat--smoke relay
affected 28\% of affected station-days at 25\,$\mu$g/m$^3$ and 19\% at 50, against
0--2\% in the co-located grid cells --- \textbf{the 10\,km analysis smooths away the peaks
that drive the relay}, and reproduces the episode at daily $r = 0.83$ while
underestimating peaks by a median factor of $\sim$2.2. Preprint, single author, single
episode; the self-directed error analysis is unusually thorough for the format.}
{10.21203/rs.3.rs-10757961/v1}

\paperentry{Qing et al.\ 2026}
{Seasonal and urban-suburban characteristics, emission and spatial sources of PM$_{2.5}$-bound PAHs and n-alkanes in Wuhu City}
{Atmos Pollut Res}
{\textbf{Metadata only.} Crossref sweep, no abstract deposited. A source-apportionment
study of the particle-bound organic fraction, which is the composition class the previous
issue's mortality result pointed at --- flagged for retry rather than characterised here.}
{10.1016/j.apr.2026.103186}

\paperentry{Liu et al.\ 2026}
{Identification of dominant weather patterns and circulation regulation mechanisms of winter PM$_{2.5}$ pollution in Iran}
{Atmos Res}
{\textbf{Metadata only.} Crossref sweep, no abstract deposited. Synoptic classification
of winter PM$_{2.5}$ episodes; one of very few Iranian PM records the watch has surfaced,
which is why it ships as a stub rather than being dropped.}
{10.1016/j.atmosres.2026.109287}

\paperentry{Abadi et al.\ 2026}
{Future projections of surface dust concentrations in major dust source regions of the Mesopotamian plains}
{Atmos Res}
{\textbf{Metadata only.} Crossref sweep, no abstract deposited. Projected mineral-dust
concentrations for a source region that exports coarse PM across the Gulf and into South
Asia; no scenario, horizon or model family is extractable from the deposited metadata.}
{10.1016/j.atmosres.2026.109282}

\paperentry{Cai et al.\ 2026}
{Effect of ammonia on reactive oxygen species generation in secondary organic aerosols}
{Atmos Environ}
{\textbf{Metadata only.} Crossref sweep, no abstract deposited, and no study location
stated. Ammonia-modulated oxidative potential of secondary organic aerosol sits precisely
between the previous issue's Nanjing ammonia-thermodynamics result and today's
oxidative-priming mechanism papers, so it is the second-highest-priority retry after
Zeng et al.}
{10.1016/j.atmosenv.2026.122307}

\vspace{4mm}
% ==================================================================== TABLE
\section*{\sffamily\bfseries\color{Deep}Corpus register}
\vspace{-1.5mm}{\color{Amber}\rule{\linewidth}{1.1pt}}\vspace{2.5mm}

{\fontsize{6.7}{8.4}\selectfont
\setlength{\tabcolsep}{2.4pt}
\begin{longtable}{@{}>{\RaggedRight}p{25mm}>{\RaggedRight}p{41mm}>{\RaggedRight}p{27mm}%
>{\RaggedRight}p{23mm}>{\RaggedRight}p{18mm}>{\RaggedRight}p{15mm}@{}}
\toprule
\textsf{\textbf{First author}} & \textsf{\textbf{Focus}} & \textsf{\textbf{Journal}} &
\textsf{\textbf{Design}} & \textsf{\textbf{Metric}} & \textsf{\textbf{Region}}\\
\midrule
\endfirsthead
\toprule
\textsf{\textbf{First author}} & \textsf{\textbf{Focus}} & \textsf{\textbf{Journal}} &
\textsf{\textbf{Design}} & \textsf{\textbf{Metric}} & \textsf{\textbf{Region}}\\
\midrule
\endhead
\multicolumn{6}{@{}l@{}}{\cellcolor{Mist}\textsf{\textbf{\textcolor{Deep}{Neuro / mental health}}}}\\[0.6mm]
\href{https://doi.org/10.1001/jamanetworkopen.2026.25015}{Dales} & Wildfire-sourced PM$_{2.5}$ and ED visits for migraine and other primary headache & \textit{JAMA Netw Open} & Case-crossover & PM$_{2.5}$ (wildfire) & Canada\\
\href{https://doi.org/10.1093/ageing/afag215}{Glover} & Umbrella meta-review of urban environmental exposures and cognitive health & \textit{Age Ageing} & Umbrella review / meta-analysis & PM$_{2.5}$ & Global\\
\href{https://doi.org/10.1016/j.nsa.2026.107019}{Goossens} & In utero air pollution exposure and pre-/postnatal brain development & \textit{Neurosci Appl} & Narrative review & PM$_{2.5}$ / PM$_{10}$ / PM$_{\mathrm{coarse}}$ & Global\\
\href{https://doi.org/10.1016/j.isci.2026.116617}{Ma} & PM$_{2.5}$-induced neurotoxicity and depression: exposure risks and mechanisms & \textit{iScience} & Narrative review & PM$_{2.5}$ + components & Global\\
\href{https://doi.org/10.3760/cma.j.cn112150-20250804-00752}{Pan} & Maternal PM exposure and offspring neurodevelopment: effects and mechanisms & \textit{Zhonghua Yu Fang Yi Xue Za Zhi} & Narrative review & PM (bulk) & China\\
\addlinespace[0.9mm]
\multicolumn{6}{@{}l@{}}{\cellcolor{Mist}\textsf{\textbf{\textcolor{Teal}{Cardiovascular \& metabolic}}}}\\[0.6mm]
\href{https://doi.org/10.3760/cma.j.cn112338-20251120-00843}{Huang} & PM$_{2.5}$ and dyslipidaemia in rural Chinese children; dietary effect modification & \textit{Zhonghua Liu Xing Bing Xue Za Zhi} & Cross-sectional (multistage) & PM$_{2.5}$ & China\\
\href{https://doi.org/10.32394/pe/218966}{Kazmierska} & Air pollution and cardiovascular disease: review of current data & \textit{Przegl Epidemiol} & Narrative review & PM$_{2.5}$ / PM$_{10}$ & Global\\
\href{https://doi.org/10.1007/s00125-026-06814-2}{Xi} & Air pollution, built environment and temperature vs. prediabetes phenotypes & \textit{Diabetologia} & Prospective cohort & PM$_{2.5}$ & Germany\\
\href{https://doi.org/10.1161/JAHA.125.046917}{Yang} & Short-term PM$_{2.5}$ and hospitalisation/ED use in atrial fibrillation & \textit{J Am Heart Assoc} & Case-crossover & PM$_{2.5}$ (wildfire-influenced) & USA\\
\addlinespace[0.9mm]
\multicolumn{6}{@{}l@{}}{\cellcolor{Mist}\textsf{\textbf{\textcolor{Sky}{Reproductive \& developmental}}}}\\[0.6mm]
\href{https://doi.org/10.1016/j.jes.2025.10.043}{Chen} & Prenatal air pollution, fetal growth trajectory phenotypes and child neurodevelopment & \textit{J Environ Sci (China)} & Prospective cohort & PM$_{2.5}$ / PM$_{10}$ & China\\
\href{https://doi.org/10.1016/j.ecoenv.2026.120566}{Huang} & Joint prenatal bisphenol + air-pollutant mixture exposure and preschool IQ & \textit{Ecotoxicol Environ Saf} & Prospective cohort (mixtures) & PM$_{2.5}$ / PM$_{10}$ & China\\
\href{https://doi.org/10.1007/s40572-026-00552-8}{Peterson} & Ambient air pollution, PAHs and pubertal development: critical review & \textit{Curr Environ Health Rep} & Critical review & PM$_{2.5}$ components & Global\\
\href{https://doi.org/10.1016/j.envres.2026.125313}{Xie} & Long-term PM$_{2.5}$ and timing of pubertal development in boys and girls & \textit{Environ Res} & Prospective cohort & PM$_{2.5}$ & China\\
\addlinespace[0.9mm]
\multicolumn{6}{@{}l@{}}{\cellcolor{Mist}\textsf{\textbf{\textcolor{Sage}{Respiratory \& allergic}}}}\\[0.6mm]
\href{https://doi.org/10.1007/s00484-026-03289-z}{Bayar Muluk} & Climate change and otolaryngology: airway disease impacts and emission reduction & \textit{Int J Biometeorol} & Narrative review & PM$_{2.5}$ & Global\\
\href{https://doi.org/10.4103/mgr.MEDGASRES-D-25-00164}{Liao} & Short-term air pollution, respiratory symptoms and lung function in asthma & \textit{Med Gas Res} & Panel study & PM$_{2.5}$ / PM$_{10}$ & China\\
\href{https://doi.org/10.3390/epidemiologia7040097}{Sanchez de Prada} & Ambient PM and asthma-coded admissions in children 0-3 y, Spain & \textit{Epidemiologia} & Ecological panel & PM$_{2.5}$ / PM$_{10}$ & Spain\\
\href{https://doi.org/10.1016/j.jhazmat.2026.143076}{Yang} & Allergenic vs. toxic risk of airborne seed hairs compared with ambient PM$_{2.5}$ & \textit{J Hazard Mater} & Field measurement & PM$_{2.5}$ + bioaerosol & China\\
\addlinespace[0.9mm]
\multicolumn{6}{@{}l@{}}{\cellcolor{Mist}\textsf{\textbf{\textcolor{Violet}{Mechanistic toxicology}}}}\\[0.6mm]
\href{https://doi.org/10.1016/j.jes.2025.12.074}{Ding} & PM$_{2.5}$ promotes vascular calcification via miR-27b-5p/HES1 signalling & \textit{J Environ Sci (China)} & In vitro + organoid + murine & PM$_{2.5}$ & China\\
\href{https://doi.org/10.1007/s00360-026-01698-5}{Maraschi} & Osmoregulatory disruption and metal bioaccumulation from settleable atmospheric PM & \textit{J Comp Physiol B} & Controlled exposure (ecotox) & Settleable APM & Brazil\\
\href{https://doi.org/10.1016/j.jes.2025.11.018}{Wu} & Macrophage NLRP3 inflammasome in real-ambient PM-induced hepatic glucose dysregulation & \textit{J Environ Sci (China)} & Real-ambient murine + co-culture & Real-ambient PM & China\\
\addlinespace[0.9mm]
\multicolumn{6}{@{}l@{}}{\cellcolor{Mist}\textsf{\textbf{\textcolor{Clay}{Occupational \& indoor}}}}\\[0.6mm]
\href{https://doi.org/10.3389/fmedt.2026.1796828}{Chen} & Nested suction instrument to cut intraoperative surgical-smoke PM$_{2.5}$ & \textit{Front Med Technol} & Bench / device experiment & PM$_{2.5}$ (surgical smoke) & China\\
\href{https://doi.org/10.1093/annweh/wxag055}{Cole} & Determinants of respirable crystalline silica exposure in Australian tunnelling & \textit{Ann Work Expo Health} & Observational exposure study & Respirable dust / RCS & Australia\\
\href{https://doi.org/10.1016/j.jhazmat.2026.143081}{Ding} & Stage-resolved hair Zn-isotope dynamics as a firefighting combustion-exposure tracer & \textit{J Hazard Mater} & Repeated-measures biomarker & Combustion-derived PM & China\\
\href{https://doi.org/10.3390/jox16040119}{Ramsden} & Xenobiotic hazards in aircraft cabin air & \textit{J Xenobiot} & Review + risk economics & Ultrafine particles & Global\\
\addlinespace[0.9mm]
\multicolumn{6}{@{}l@{}}{\cellcolor{Mist}\textsf{\textbf{\textcolor{Amber}{Sensing, forecasting \& instrumentation}}}}\\[0.6mm]
\href{https://doi.org/10.1016/j.jenvman.2026.130515}{Shi} & Adaptive spatiotemporal graph transformer for multi-site multi-step PM$_{2.5}$ forecasting & \textit{J Environ Manage} & Deep-learning model & PM$_{2.5}$ & China / India\\
\href{https://doi.org/10.3390/bios16070380}{Tsuchiya} & Seasonal and indoor-outdoor characteristics of fluorescent aerosol particles in offices & \textit{Biosensors} & Real-time sensor campaign & Bioaerosol / coarse & Japan\\
\addlinespace[0.9mm]
\multicolumn{6}{@{}l@{}}{\cellcolor{Mist}\textsf{\textbf{\textcolor{Amber}{Exposure assessment \& modelling}}}}\\[0.6mm]
\href{https://doi.org/10.1029/2025GH001621}{Solorzano-Araque} & Multiscale exposure modelling: LOTOS-EUROS + mobility surveys + particle toxicity & \textit{GeoHealth} & CTM + mobility integration & PM$_{2.5}$ / PM$_{10}$ & Colombia\\
\addlinespace[0.9mm]
\multicolumn{6}{@{}l@{}}{\cellcolor{Mist}\textsf{\textbf{\textcolor{Coral}{Burden, policy \& mitigation}}}}\\[0.6mm]
\href{https://doi.org/10.1002/hsr2.72886}{Karamzad} & Ischaemic heart disease burden in MENA 1990-2021; PM attributable fraction 28.2\% & \textit{Health Sci Rep} & GBD secondary analysis & Ambient PM & MENA\\
\href{https://doi.org/10.1371/journal.pone.0347989}{Rees} & HEPA filtration to reduce respiratory infection in care homes & \textit{PLoS One} & RCT process evaluation & Filtered PM & UK\\
\href{https://doi.org/10.1016/j.jes.2026.01.009}{Shan} & Policy-driven emission reductions and post-pandemic rebound, high-resolution inventories & \textit{J Environ Sci (China)} & Emission inventory & PM emissions & China\\
\addlinespace[0.9mm]
\multicolumn{6}{@{}l@{}}{\cellcolor{Mist}\textsf{\textbf{\textcolor{Slate}{Other clinical endpoints}}}}\\[0.6mm]
\href{https://doi.org/10.3390/dj14070440}{Lupita} & PM$_{2.5}$, preventive-care capacity and caries/edentulism burden across the EU & \textit{Dent J} & Ecological panel & PM$_{2.5}$ & EU-27\\
\href{https://doi.org/10.1016/j.jdin.2026.06.008}{Park} & Long-term PM exposure and risk of androgenetic alopecia & \textit{JAAD Int} & Nationwide cohort & PM$_{2.5}$ / PM$_{10}$ & South Korea\\
\href{https://doi.org/10.3171/2026.3.PEDS2615}{Sun} & Environmental exposure and area deprivation vs. paediatric hydrocephalus outcomes & \textit{J Neurosurg Pediatr} & Registry cohort + geospatial & PM$_{2.5}$ (satellite-derived) & USA\\
\addlinespace[0.9mm]
\bottomrule
\end{longtable}}


\vspace{4mm}
\begin{keybox}
{\sffamily\bfseries\footnotesize\color{Deep}Provenance and method}\\[1.2mm]
{\fontsize{7.2}{8.8}\selectfont
Entry window \textbf{2026-08-22 $\rightarrow$ 2026-08-22}, contiguous with the 21 August
issue. Built at 22:4x local on 22 August, satisfying the 22:00 rule for same-day
harvesting. \textbf{Sources.} PubMed E-utilities on \texttt{[EDAT]} returned only
\textbf{2 records}, both shipped --- a Saturday floor. Europe PMC on
\texttt{CREATION\_DATE} returned 23 unique records, of which \textbf{15 were archival PMC
back-scans} (legacy scans, no title, publication year decades outside the window), 3 were
already seen or belonged to the 21 August issue, and 3 shipped. The Crossref by-ISSN
sweep returned 36 records after front-matter filtering and \textbf{shipped 8 --- 62\% of
the issue}. \textbf{OpenAlex returned HTTP 429 for the 20th consecutive run.} arXiv
returned nothing new under the 3-day screen.
\textbf{Two harvester changes this run.} The supplementary ISSN list trialled on 20 August
was promoted into \texttt{harvest.JOURNALS} (15 additions, each ISSN verified against
Crossref's journal index; MDPI \textit{Sensors} deliberately excluded after returning 27
deposits for 0 in-scope records). The first run with \textit{Environ Sci Technol} in the
list returned 154 works of which \textbf{142 were mastheads and issue-information stubs}
being re-deposited in bulk from volumes 42--52; a title-pattern front-matter filter now
drops these at the source, and the count is logged per journal.
\textbf{Screening.} 44 records were screened out and written to
\texttt{state/rejected.jsonl} with an individual reason each. Two borderline exclusions
are logged as borderline: combustion SVOC residue on wildland-firefighter protective gear,
and droplet-residue resuspension from medical PPE --- both occupational and
combustion-adjacent, neither carrying an airborne particulate measurement.
\textbf{Labels.} Six records carry no abstract and are marked \textit{Metadata only};
their design, sample size, endpoint and (in two cases) study location are recorded as not
extracted rather than inferred. Effect estimates are transcribed verbatim from abstracts.
All DOIs were checked by \texttt{check\_dois.py} before publication. Content derived from
publisher metadata; abstracts remain the copyright of their respective publishers.\par}
\end{keybox}

\end{document}
